\section{Equal-chunk interchange from the finite rational-frame motif}
\label{app:swap}

The finite rational-frame network and array-interchange strategy are
adapted from OpenAI's Sections 3--4~\cite{swap}. The explicit reconstruction
below supplies the parameters and tape accounting used in this paper.

This appendix constructs, rather than assumes, a fixed-tape interchange
algorithm.  For every binary width $u\ge1$, positive integers $P,G,B$,
and arbitrary $B$-bit record payload, the machine takes an array in
lexicographic order on
\[
 [P]\times[2^u]\times[G]\times[2^u]\times[B]
\]
to the array whose two $u$-bit fields are exchanged, retaining $P,G,B$
and the internal order of each record.  Its time is
\[
                        O(Vu^\tau),
\]
where $V=PGB2^{2u}$ is the entire stored volume and one fixed exponent
$\tau=1-2^{-36}<1$ suffices.  All counters, fixed-table preparation,
copying, row padding, recursive returns, and cleanup are included.  No
unverified global multiplication algorithm or equal-chunk interchange
oracle is a premise.

\subsection{The finite scalar routing network}

Let $\mathcal T=\binom{[g]}3$ with $g=48$.  Two triples are adjacent when
their intersection has size one.  Put
\[
 v=\binom{48}{3}=17296,\quad z=3\binom{45}{2}=2970,
 \quad N=v^3,\quad I=3v^2.
\]
For $S,T\in\mathcal T$, define the binary maps
\[
 \begin{aligned}
 V_{(S,T),T'}&=[T=T']\,[|S\cap T|=1],\\
 J_{S,(S,T)}&=[|S\cap T|=1],\\
 G_{i,T}&=[i\in T],\qquad R_{S,i}=[i\in S].
 \end{aligned}
\]
Then $(JV)_{S,T}$ is one exactly when $|S\cap T|=1$, and
$(RG)_{S,T}=|S\cap T|\pmod2$.  If $S=T$, the intersection size is
three and only $RG$ contributes one.  For distinct triples with
intersection one both terms contribute and cancel in $\mathbb F_2$;
for intersection zero or two both are zero.  Hence
\[
                              RG+JV=I_v.
\]

One invocation has source $x$, target $y$, arbitrary side scratch $a$,
and arbitrary central scratch $c$.  The chronological eight rows are
\[
\begin{aligned}
 y&\gets y-Ja,& y&\gets y-Rc,& a&\gets a+Vx,& c&\gets c+Gx,\\
 y&\gets y+Rc,& y&\gets y+Ja,& c&\gets c-Gx,& a&\gets a-Vx.
\end{aligned}
\]
For initial scratch $(a_0,c_0)$ the target increment is
\[
 -Ja_0-Rc_0+R(c_0+Gx)+J(a_0+Vx)=(RG+JV)x=x,
\]
and the final scratch is exactly $(a_0,c_0)$.  The signs are harmless
over $\mathbb F_2$, but writing them makes the restoration identity
visible.  Thus the invocation adds its source to its target for every
initial scratch value.

Use three stages.  The first is forward from bank $X$ to bank $Y$, the
second is the inverse schedule from $Y$ to $X$, and the third is forward
from $X$ to $Y$.  On arbitrary bank contents,
\[
 (x,y)\longmapsto(x,y+x)\longmapsto(-y,y+x)\longmapsto(-y,x).
\]
Over $\mathbb F_2$ this exchanges $X,Y$.  Each triple coordinate
participates in each stage once; invocations have disjoint scratch
roles.  All scratch roles are restored.  The resulting finite role
permutation $\rho$ exchanges the $N$ data roles and fixes every scratch
role.  It is the scalar-routing network used below.  The finite incidence
identity and scratch schedule are also recorded in
\path{supporting/SCALAR_MOTIF_SYMBOLIC_AUDIT.md} and
\texttt{finite\_phase\_appendix.tex}.

\subsection{Rational frames and exact child count}

Address arithmetic is over $\mathbb Q$ first and then over
$\mathbb Z/(q^b)$ for a fixed odd prime $q$.  It is distinct from the
payload network over $\mathbb F_2$.  On $F=\mathbb Q^{48}$ use the
nondegenerate symmetric form with Gram matrix
\[
                         B_0=I-J/9.
\]
The all-one line has eigenvalue $1-48/9=-13/3$, and its orthogonal
complement has eigenvalue one, so the form is nondegenerate.  For
$U\in\mathcal T$, let $t_U=\mathbf1_U$.  Then
\[
 \langle t_U,t_V\rangle=|U\cap V|-1,\qquad
 \langle t_U,t_U\rangle=2.
\]
Thus $P_U$, the orthogonal projection onto $\langle t_U\rangle$, is a
well-defined rational rank-one projection.  Adjacent triples give
orthogonal lines.  The common-label table for the finite motif follows
from these orthogonal direct-sum decompositions; a derivation is in
  \path{supporting/ntt_base_five/finite_motif_rank_closed_form_audit.md}.

At one frame stage, let $A$ be the preceding tensor label, $P$ its
prefix line, $B=P^\perp$, and $L_U=\langle t_U\rangle$.  The next center
label satisfies
\[
 (P\otimes L_U)^\perp
   =(B\otimes F)\oplus^\perp(P\otimes L_U^\perp).
\]
Side inclusions use $L_U\perp L_V$ for adjacent $U,V$.  In the fixed
physical gate order the common subspace labels are
\[
\begin{array}{c|l}
0&B\otimes L_U\\
1&B\otimes F\\
2&(B\otimes F)\oplus^\perp(P\otimes L_U)\\
3&A\otimes F\\
4&B\otimes F\\
5&(B\otimes F)\oplus^\perp(P\otimes L_U^\perp)\\
6&A\otimes F\\
7&A\otimes F.
\end{array}
\]
Every label is tensored with the fixed suffix tensor line of the
invocation. A side wire named by its physical Y/X neighbor pair $(S,T)$
meets rows $0(S),2(T),5(S),7(T)$; its middle inclusion is
$L_T\subset L_S^\perp$. A center wire meets rows $1,3,4,6$.
Data X meets $2,3,6,7$ and data Y meets $0,1,4,5$.
The second stage uses the reversed schedule with logical banks exchanged,
so its ordered side-pair names are reversed to retain this physical
convention. Scratch input frames are zero, scratch output frames are
$I_m$, and the final data X/Y frames are $I_m$ and $I_m-P_U$.
These finite inclusions and wire incidences use no phase-kernel,
weighted-mask, H2, or interchange algorithm as a premise.

In the fixed
eight-role order the label dimensions are
\[
 a-1,\quad 48(a-1),\quad48(a-1)+1,\quad48a,\quad
 48(a-1),\quad48a-1,\quad48a,\quad48a,
\]
where $a=\dim A$.  The only decreases are the central transitions
$48a\to48(a-1)$, each losing 48 dimensions.  With the fixed counts
\[
 W=2N+I(vz+48),\qquad m=48^3=110592,
 \qquad L=I\cdot48^2,
\]
the signed endpoint dimension telescope and twice the downward variation
give the sum of comparable-edge ranks $Wm-2N+2L$.

\paragraph{Complete physical edge ledger.}
The inverse middle stage runs logical rows $7,6,\ldots,0$ with source
Y and target X. Its physical gate incidences are nevertheless
Y-side, Y-center, X-side, X-center, Y-center, Y-side, X-center, X-side,
the same chronological order as the other stages. Thus the common-label
table is used in its printed order, with the side-pair names reversed;
one must not reverse the table along with the logical row numbers.
Put $d=48$ and $a=d^{j-1}$ at stage $j$. Suppressing the suffix line,
the complete scratch dimension paths, including their endpoints, are
\[
\begin{array}{c|l|c}
 \hbox{physical wire}&\hbox{dimension path}&\hbox{downward loss}\\ \hline
 \hbox{side}&0,\ a-1,\ d(a-1)+1,\ da-1,\ da,\ m&0\\
 \hbox{center}&0,\ d(a-1),\ da,\ d(a-1),\ da,\ m&d .
\end{array}
\]
These are geometric inclusions: the middle side edge uses
$L_T\subset L_S^\perp$, whose nondegenerate residual has dimension
$d-2$; the center's downward residual is $P\otimes F$.
At an interstage boundary X enters as $A\otimes L_T$, contained in
$(B\otimes F)\oplus(P\otimes L_T)$, while Y enters as
$B\otimes L_S$, exactly its next first label. The physical X path is
therefore nondecreasing from a formal positive input line to dimension
$m$, and Y is nondecreasing from zero to dimension $m-1$.
Their ordinary edge-rank sums are both $m-1$.
Each side contributes $m$ and each center $m+2d$.
Since there are $I\cdot48$ centers, summing gives directly
\[
 2N(m-1)+(W-2N)m+2I\cdot48d=Wm-2N+2L.
\]
The actual negative X source, treated next, replaces its formal
zero-rank first edge by one rank-one edge; it does not alter the
ordinary interstage inclusions.

At each of the $N$ exceptional $X$ sources choose input frame $-P_U$;
the first gate frame is $+P_U$.  Their difference is
$2P_U$, of rational rank one.  Since the address modulus will have odd
characteristic, this is not zero.  These $N$ source corrections add $N$
recursive children, so the full child count is
\[
                         s=Wm-N+2L.
\]
For every routed role $w$, the endpoint labels obey
\[
                  M_{\rm out}(\rho(w))-M_{\rm in}(w)=I_m.
\]
For an $X$ data role the endpoint difference is
$P_{U^\perp}-(-P_U)=I_m$; the other data and scratch endpoint pairs
are the corresponding zero/full or complementary-label pairs in the
finite common-label table.  For comparable labels, the rank of the
projection difference is the dimension change.  The only equal-dimension
exception is the $N$ first $X$ edges, already counted as $2P_U$ of rank
one.  Thus no source correction is hidden by the endpoint identity.
At each role edge with label difference $A$, implement the corresponding
address-frame change.  The scalar routing sends role $w$ to $\rho(w)$;
the displayed endpoint identity then leaves the same target shear on
data and scratch roles alike.

The exact fixed counts reduce to
\[
 \frac WN=\frac{9633881}{1081},\qquad
 \frac LN=\frac{432}{1081},\qquad
 1-\frac{s}{Wm}=\frac{217}{9633881\cdot110592}.
\]
The integer inequality
\[
 217\cdot2^{36}>12\cdot9633881\cdot110592
\]
and $m<2^{17}$ give $\ln m<12$ and
$1-s/(Wm)>12\cdot2^{-36}$.  From $e^{-x}>1-x$ it follows that
\[
                      s/W<m^\tau,\qquad \tau=1-2^{-36}.
\]
All constants here are fixed.  These are exact rational counts, not a
finite-size matrix-rank inference.

\subsection{Choosing one odd address prime}

There are finitely many rational label matrices, projections, and frame
differences.  For each edge matrix the construction below gives a fixed
factorization into lower-triangular matrices and a partial permutation.
Collect the denominators of all entries and the numerators and
denominators of every nonzero diagonal unit in these factorizations and
their inverses.  This is a finite set of nonzero integers.  Choose one
odd prime $q$ dividing none of them.  Such a prime exists by Euclid's
argument; trial division through primes finds one after finitely many
steps.  This is fixed finite setup, independent of $b,u,P,G,B$.  The
prime and the rational matrices may equivalently be constants of the
finite machine.

All rational matrices now reduce to $\mathbb Z/(q^b)$ for every $b$.
The chosen triangular diagonal entries remain units.  Each fixed
factorization remains an exact identity, with the same $a$ partial
permutation pivots.  No data-dependent division or prime oracle is used.
Because $q$ is odd, the exceptional source matrix $2P_U$ still has its
one nonzero pivot and costs one child swap.

\subsection{Fixed-radix tape operations}

The following operations use fixed $q$ and fixed tape families.

\paragraph{Counters and records.}
Use little-endian ripple counters with the head returning to the low
end. Over a run of $A$ increments, or a run of $A$ decrements, of an $r$-bit
counter from an arbitrary initial state, bit $j$ is visited
$O(1+A/2^j)$ times, so the total cost is $O(A+r)$.
The width term is necessary: one increment from $2^r-1$ can visit
all $r$ bits. For a nonempty range of $n$ values, use its own
$r=\max\{1,\lceil\log_2(n+1)\rceil\}$ active counter width and
initialize/reset it at the prescribed endpoint. A full scan,
initialization and return then cost $O(n+r)=O(n)$.
To test termination without a full-width comparison after each step,
initialize a countdown to $n-1$, execute the loop body, then
decrement and stop precisely on borrow out of its marked $r$-bit
interval. The last decrement from zero performs the end test
while visiting those bits. This gives exactly $n$ bodies, including
$n=1$. A needed forward index uses a separate amortized counter.
Range bounds are prepared once in their own canonical widths;
each reset copies only that bound, not a full ambient-width register.
Stale cells outside this active interval are not rescanned at every
reset. With a fixed number of nested ranges, charge each inner reset
to its ensuing nonempty scan. Copying a $B$-bit record thus costs
$O(B)$ including its counter initialization and returns, even for
$B=1$. Recursion saves counters/descriptors at the current stack top;
it does not cross an ancestor to fetch a length. Ordinary schoolbook
shape arithmetic has a fixed polynomial cost in descriptor bitlength;
the retained-range bound below pays it at every recursive call.

More explicitly, traversing a fixed-depth product of nonempty ranges
$n_1,\ldots,n_d$, including a final bit-record range if necessary,
has counter cost bounded by
\[
 O\!\left(\sum_{i=1}^d
       \Bigl(\prod_{j<i}n_j\Bigr)
       (n_i+\log(2n_i))\right)
   =O\!\left(d\prod_{i=1}^d n_i\right).
\]
For the current fixed matrix dimension, $d$ is fixed, even when
prefix, gap and suffix ranges are arbitrarily large. All array
bits remain raw payload: counters and boundary markers are on
work tapes, with no growing header added to each record. Rewinding
the fixed number of current split tracks costs their combined
volume and stops at their current stack boundaries.

\paragraph{Digit moves and pointwise maps.}
Split a stream into $q$ tracks by one selected radix-$q$ digit, preserving
order on each track, and merge in the desired digit order.  A single
digit move costs $O(V)$.  For a fixed permutation of two digit positions,
split into $q^2$ tracks and merge according to the two preimage digits;
this also costs $O(V)$.  A fixed pointwise finite-alphabet map on aligned
tracks costs $O(V)$.  The number $q^2$ is large but fixed, not input
dependent.

\paragraph{Rotations controlled by earlier fields.}
Let the target field have $Q=q^b$ values and follow every field used to
compute its shift.  Within each fixed earlier-address fiber, it consists
of $Q$ consecutive blocks of suffix records.  For shift $a\in[Q]$,
copy the first $Q-a$ blocks to one stream and the last $a$ to another,
then write the latter before the former.  This realizes
$y\mapsto y+a\pmod Q$ in $O(V)$ record motion.  Computing the offset by
schoolbook arithmetic on the fixed number of $b$-digit controls costs
$\operatorname{poly}(b)$ per fiber.  Each fiber contains $q^b$ blocks,
so this metadata is absorbed into its own volume.  Arbitrary spectators
and payload remain inside each block.
The enumeration retains separate counters for the selected controls.
Large spectator counters are advanced or reset by the preceding
product-scan rule; their full values are not decoded or copied into
the offset calculation. Thus the bound is $O(V)$ even for $b=1$,
with arbitrarily large prefix, intervening gap, or suffix.
This argument relies on reading only the fixed number of selected
$b$-digit fields. It does not authorize an offset depending on the
entire ambient address at the same cost.

\paragraph{Fixed rational units.}
For a fixed positive integer $c$ coprime to $q$, partition $y\in[Q]$
into its $c$ monotone pieces indexed by $j=\lfloor cy/Q\rfloor$.
On a piece the output is $y'=cy-jQ$ in increasing order.  Since
$Q$ is invertible modulo $c$, for each output value $y'$ the unique piece
is determined by $j\equiv-y'Q^{-1}\pmod c$.  Split by these fixed $c$
classes and merge their increasing streams in output order.  This is
$y\mapsto cy\pmod Q$ in $O(V)$, including piece endpoints and counters.
Prepare the endpoints and $Q^{-1}\bmod c$ in polynomial-$b$ time;
the output selector needs only $y'\bmod c$, a fixed finite-state
counter updated once per suffix block. It performs no full $b$-bit
multiplication or global-address decoding at each record.
The same residue split on a supplied output stream constructs the inverse
permutation directly.  Negation modulo $Q$ fixes zero and reverses the
other $Q-1$ blocks; a stack reversal plus reversal within each block costs
$O(V)$.  A fixed rational unit $c/d$ is the composition of integer
multiplication by $|c|$, inverse multiplication by $|d|$, and a possible
negation.  The prime $q$ avoids $c,d$.

\subsection{Lower-triangular frames and rank-many child swaps}

For a fixed rational matrix $A$ of rank $a$, eliminate its topmost
nonzero row using the rightmost nonzero pivot.  Clear entries to the
left by lower-triangular column operations and entries below by
lower-triangular row operations.  Remove the isolated pivot row and
column and continue in the remaining ordered indices.  This yields
\[
                         A=E_1\Pi E_2,
\]
where $E_1,E_2$ are nonsingular lower-triangular matrices and $\Pi$
has exactly $a$ ones in distinct rows and columns.  Reversing the
elimination proves the factorization; nonsingularity of $A$ is not
required.
Concretely, if the pivot is in row $i_0$, column $j_0$, clearing a
column $j<j_0$ adds a multiple of column $j_0$ to column $j$,
a lower-triangular right operation. Clearing a row $i>i_0$ adds a
multiple of row $i_0$ to row $i$, a lower-triangular left operation.
Normalize the pivot by a nonzero diagonal scale. Previously isolated
pivot rows and columns are unchanged, since their entries in the
remaining submatrix are zero. Retaining the original order on the
remaining indices preserves triangularity. Thus $LAR=\Pi$ with
invertible lower-triangular $L,R$, and $E_1=L^{-1},E_2=R^{-1}$.
The partial permutation need not be diagonal.

\begin{lemma}[Signed factorization over every address power]
\label{lem:swap-signed-module}
For each actual directed rational frame difference $A=M_j-M_i$ of
rank $a$, the chosen fixed prime gives, for every $b\ge1$,
\[
 A_b=E_{1,b}\Pi E_{2,b}\quad\hbox{over }R_b=\mathbb Z/(q^b),
\]
with both triangular factors invertible. The image and kernel of $A_b$
are free direct summands of ranks $a$ and $m-a$. The same $a$ unit
pivots implement the signed address shear; no field assumption on
$R_b$ or invertibility of a data coordinate is required.
\end{lemma}
\begin{proof}
The finite rational identities reduce entrywise since $q$ avoids every
entry denominator. Each triangular diagonal remains a unit by the
same fixed exclusion, so both factors are invertible over every $R_b$.
Multiplication by those invertible factors identifies the image and
kernel with those of $\Pi$, whose nonzero entries are literal ones.
In particular these are coordinate direct summands of the stated ranks.
This is an explicit factorization certificate over a ring, rather than
a rank test over a field.

Factor the \emph{signed} matrix $A$ before reduction. Its negative
can equally be factored as $(-E_1)\Pi E_2$, absorbing the sign in
unit diagonals. For the source correction $2P_U$, the factor two
is likewise absorbed in one triangular factor; $q$ is odd.
It retains one pivot and one child. No sign or factor two is reduced
using the separate characteristic-two payload arithmetic.
\end{proof}

A lower-triangular coordinate change within a bank is performed in
descending coordinate order.  At target coordinate $i$, the controls
with index $j<i$ still have their original values and physically precede
the target.  Apply its fixed diagonal unit, then add the fixed linear
combination of those earlier coordinates by later-target rotations.
This costs $O(V)$ per coordinate; the dimension $m$ is fixed.

Define, for vectors of $m$ radix-$q^b$ digits,
\[
 \Phi_A(H,D)=(H+AD,D).
\]
Its induced forward record pushforward is
\[
 (\mathcal P_A f)(H,D)=f(H-AD,D),\qquad
 \mathcal P_A\mathcal P_B=\mathcal P_{A+B},\qquad
 \mathcal P_A^{-1}=\mathcal P_{-A}.
\]
Here the coordinate signs are in $R_b$; pointwise payload addition
remains XOR over $\mathbb F_2$.
To implement it, first change $D$ by $E_2$, change $H$ by $E_1^{-1}$,
add $\Pi D$ to $H$, undo the $H$ change by $E_1$, and finally undo the
$D$ change by $E_2^{-1}$.  The result is
$H+E_1\Pi E_2D=H+AD$ with the original $D$ restored.  The triangular
operations are linear volume.
For each one in $\Pi$, the cross-bank shear $H_i\gets H_i+D_j$ is
implemented in the physical order (earlier field $H_i$, later field
$D_j$) by
\[
       D_j\gets D_j+H_i;\qquad
       \operatorname{swap}(H_i,D_j);\qquad
       D_j\gets H_i-D_j.
\]
On $(h,d)$ this is $(h,d)\mapsto(h+d,d)$.  The first and last
operations target later fields and cost $O(V)$.  The middle operation is
one recursive interchange of two $b$-digit fields across arbitrary gap
and suffix.  Thus $\Phi_A$ costs exactly $a$ recursive child swaps plus
$O(V)$ other work.  Its inverse is $\Phi_{-A}$ and has the same bound.
Lemma~\ref{lem:swap-signed-module} justifies these same fixed factors
for singular, negative and source-correction matrices at every $b$.

\subsection{Exact power widths, role frames, and tape recurrence}

First let each chunk have width $e=m^k$ radix-$q$ digits.  Keep an
independent row field of size $R$ divisible by $W^k$, before both chunks;
arbitrary prefix, gap, suffix, and spectators between the row and first
chunk are retained.  For $k=0$ the two $q$-digit chunks are exchanged by
the fixed split/merge construction in $O(V)$.

For $k>0$, split each chunk into $m$ consecutive width-$b$ fields,
$b=e/m$, and denote their vectors by $H,D$.  Split whole rows cyclically
into $W$ role streams inside each earlier prefix, stream $w$ taking rows
$r=Wa+w$.  Each role has the full $H,D$ ranges and all spectators, and
row count $R/W$ divisible by $W^{k-1}$.  On each finite scalar-network
edge from label $M_i$ to $M_j$, apply $\Phi_{M_j-M_i}$ to its role stream,
then perform the specified pointwise XOR gate on aligned payloads.

For input array $f_w$ on role $w$, use the logical interpretation
$\mathcal P_{-M_{\rm in}(w)}f_w$.  This is a proof invariant, not a
preprocessing pass.  The frame change on an edge is common to its aligned
XOR operands and commutes with pointwise addition.  The scalar network
routes role $w$ to $\rho(w)$; its physical output there is
\[
\mathcal P_{M_{\rm out}(\rho(w))-M_{\rm in}(w)}f_w
       =\mathcal P_I f_w.
\]
Equivalently, with block-diagonal input/output pushforward frames
$\mathcal D_{\rm in},\mathcal D_{\rm out}$ and scalar role
permutation $P$, the entire framed word is
\[
 \mathcal D_{\rm out}P\mathcal D_{\rm in}^{-1}
        =P(I_W\otimes\mathcal P_I).
\]
The endpoint difference is $I_m$ on every routed role, including dirty
side and center roles. A common reindexing commutes with their XOR
gates; no payload value is assumed zero.
Merge output role $\rho(w)$ back into input row class $w$.  This undoes
the role exchange and applies $H\gets H+D$ to every row, including
arbitrary scratch rows, while preserving all spectators.

The finite role network supplies only this earlier-target shear.  Complete
the chunk interchange with the outer later-target operations
\[
 D\gets D-H,\qquad H\gets H+D,\qquad D\gets H-D.
\]
On input $(h,d)$ the result is $(d,h)$.  Only the middle operation uses
the framed role network.  The first and last operations are later-target
rotations and fixed unit maps, each costing $O(V)$.  Thus the number of
recursive children remains exactly $s$; it is not multiplied by three.

The finite edge matrices have ranks summing to $s$. A pivot child is a
direct call to the exact power-width entry. More explicitly, its parent
role shape is
\[
\begin{aligned}
 &[P][R/W][A][H_1]\cdots[H_m]\\
 &[G][D_1]\cdots[D_m][B].
\end{aligned}
\]
Every formal field has width $b=e/m$, with a separate pretarget
spectator range $A$. A pivot using $H_i,D_j$ has child parameters
\[
 A_c=Aq^{b(i-1)},\qquad
 G_c=q^{b(m-i)}Gq^{b(j-1)},\qquad
 B_c=q^{b(m-j)}B.
\]
The call is
$\operatorname{SwapPower}(b,q,P,R/W,A_c,G_c,B_c)$.
The row range remains explicit and is divisible by $W^{k-1}$;
all unselected fields remain in the contiguous spectator products.
Consequently
\[
 P(R/W)A_cG_cB_cq^{2b}
   =P(R/W)AGBq^{2e}=V/W.
\]
This power entry introduces no numeric or row pad. It does not call
the arbitrary-width or binary wrapper. The pivot factorization
therefore gives exactly $s$ recursive calls of volume $V/W$ and
width $e/m$. The
finite number of role gates, fixed-rational operations, XOR scans, row
splits/merges, and setup work contributes $C_1V$ at the parent.  Use
separate fixed tape families and a parking stack: before one child, park
only the other current role streams at the stack top; pass the active
stream to the child; after return, restore siblings by popping and writing
right-to-left.  Clear only current visited intervals.  A descendant never
traverses a parked ancestor.  Since $s,W$ are fixed constants, all parent
parking costs $C_sV$ and does not multiply child time.

With normalized time $F_k$, the exact recurrence is
\[
 F_0\le C_0,\qquad
 F_k\le(s/W)F_{k-1}+C_1.
\]
By the finite count $s/W<m^\tau$, this gives
\[
                         F_k=O(m^{k\tau})=O(e^\tau).
\]
The bound also holds if $s/W\le1$, since the resulting linear sum in
$k$ is at most $O(e^\tau)$.  Parent parking uses $O(V)$ time and
storage at the current level.  At depth $j$, the child volume is
$V_j=V_0/W^j$; the two original full chunk ranges remain present as
targets or spectators, so $V_j\ge q^{2e_0}$ for root width $e_0$.
Since $j\le k=O(\log e_0)$,
\[
 \log V_0=\log V_j+j\log W=O(\log(2V_j)).
\]
Any fixed polynomial of the root descriptor bitlength therefore fits
within $O(V_j)$ at the child.  Its descriptor lists only the current
fields and collapsed spectator intervals, not the chain of ancestor
subdivisions.

\subsection{Arbitrary $q$-radix widths without rounding the core}

Let the two original chunks have width $e$ radix-$q$ digits.  Put
$k=\lceil\log_m e\rceil$.  Choose the least $r_0$ such that
$q^{2r_0}\ge W^k$.  For sufficiently large $e$, $1\le r_0<e$ and
$r_0=O(\log(2e))$; the bounded smaller range uses direct finite
single-digit interchanges and changes only the uniform constant.
Write the first chunk as $h_+h_-$ and the second as $d_+d_-$,
with the four lengths specified separately:
\[
 |h_+|=|d_+|=r_0,\qquad |h_-|=|d_-|=e-r_0.
\]
Thus the plus sign always means high digits and the minus sign
always means low digits. Starting from
\[
 [h_+][h_-][G][d_+][d_-],
\]
first exchange corresponding high digits individually, obtaining
\[
 [d_+][h_-][G][h_+][d_-].
\]
Then move only the $r_0$ digits of $h_+$ left across $h_-$ and $G$
using fixed-radix digit moves, obtaining
\[
 [d_+][h_+][h_-][G][d_-].
\]
This costs $O(Vr_0)$.  The joined row field $(d_+,h_+)$ has size
$R=q^{2r_0}\ge W^k$.  Pad this one independent numeric row range to
\[
 R'=W^k\left\lceil\frac{R}{W^k}\right\rceil,
                   \qquad R\le R'<2R.
\]
It is divisible by every $W^j$ needed in the recursion.  This private
factor below two multiplies the current volume by one constant factor for
the whole arbitrary-width call.  All children receive that same padded
row range and have volume $V'/W$ at the first level (then divide by $W$
per level); the pad is not independently reintroduced at descendant
levels.  Pad rows are initialized to zero.  Completed interchange calls
preserve the row coordinate, so these rows remain zero at the crop
boundary.  Crop the row factor once, after the complete arbitrary-width
core, to recover the original volume.

Write $e-r_0$ in base $m$.  At scale $m^j$ there are at most $m-1$
matching consecutive pieces in each lower chunk.  Interchange each
corresponding pair by the exact power-width routine, with all other
pieces as spectators.  The sum of their costs is
\[
 O(V)\sum_j a_jm^{j\tau}=O(Ve^\tau),\qquad 0\le a_j<m.
\]
No target piece is rounded up to a larger power.  After all pieces,
the field order is
\[
 [d_+][h_+][d_-][G][h_-].
\]
Crop the row padding back to $R$ genuine rows before undoing the high
digit move.  Move $h_+$ after $G$ by the reverse-direction sequence of
$r_0$ digit moves, obtaining
\[
 [d_+][d_-][G][h_+][h_-].
\]
The two original chunks have been interchanged.  The high moves cost
$O(V\log(2e))=O(Ve^\tau)$, and row padding contributes only a fixed
factor below two.  All spectators retain their order.

For completeness, fixed-radix digit moves use $q$-way stream splitting
and merging; their cost is $O(V)$ each and uses a fixed number of tapes.
The $O(\log e)$ piece schedule and shape arithmetic are polynomial in
$\log(2V)$, hence $O(V)$.  The full chunk ranges remain in every
descendant, which is the volume lower bound used for descriptor costs.

Counters and record movement use ordinary sequential tapes.  A
little-endian binary counter changes bit $j$ only $O(1+T/2^j)$ times
over $T$ consecutive increments, so a run of $T$ increments, its reset,
and return cost $O(T+w)$ for width $w$. Use a range-specific active
width as above; do not restart a full ambient-address counter for
each constant-size inner scan. Nested counters and child
descriptors are pushed at the current stack top and only that child's
interval is popped on return.  A $B$-bit suffix record is copied with an
amortized down-counter in $O(B)$, including reset.  At each call there
are only a fixed number of shape parameters; polynomial work in
$\log(2V)$ is $O(V)$ for all positive input volumes after increasing the
uniform constant.  No counter repeatedly scans an ancestor descriptor.

\subsection{Binary fields and final theorem}

For binary width $u$, set $N=2^u$ and choose the least $e$ with
$Q=q^e\ge N$.  Then $Q<qN$.  Pad each of the two numeric address ranges
from $N$ to $Q$, copying the valid rectangle and filling its complement
with zero records.  This nested scan costs $O(V')$ where
\[
 V'=PGBQ^2<q^2V.
\]
Radix-$q$ lexicographic order is numeric order, so no per-record base
conversion is used.  Apply the exact $q$-digit interchange above.  Its
completed coordinate transposition preserves the nested valid rectangle
$[N]\times[N]$; hence invalid padding remains zero and can be cropped
after the full operation.  No crop is performed inside a recursive piece.
Because $e=O(u)$ for fixed $q$, the final time is
\[
                         O(V'u^\tau)=O(Vu^\tau).
\]

All operation families have fixed finite tape sets: the $q^2$ split
tracks, the $W$ role streams, reusable I/O/work tapes, and stack/descriptor
tapes depend only on fixed constants.  Finite rational tables and the
avoiding prime are fixed machine data or are generated by terminating
constant setup.  Counters use ordinary amortized binary increments;
record copies and returns cost their full bit volume.  Thus this is one
uniform fixed-finite-tape construction for all $P,G,B,u$, including
$B=1$.  It proves the local interchange statement without assuming the
source preprint's global integer multiplication algorithm.  It makes no
novelty claim for the finite motif or its routing method.



